\documentclass[12pt]{article}
\usepackage[T1]{fontenc}
\usepackage[utf8]{inputenc}
\usepackage{lmodern}
\usepackage{textcomp}
\usepackage{enumitem}
\usepackage{geometry}
\geometry{margin=1in}
\begin{document}
\begin{center}
Answer Key - Economics - Undergraduate Level Assessment 6 \
\small (Answers are ordered exactly as in the question paper.)
\end{center}

\section*{Multiple Choice}
\begin{enumerate}[label=\arabic*.]
\item \textbf{Answer:} A
\\ \textit{Options:} A) Geographic diversity of farming regions; B) Environmental concerns; C) Consumer preference for organic produce; D) Lack of funding; E) Technological advancements in agriculture
\\ \textit{Note:} The geographic diversity of farming regions means that agricultural practices, challenges, and economic conditions vary widely across locations, making it difficult to implement a one-size-fits-all income distribution policy for farmers.
\item \textbf{Answer:} A
\\ \textit{Options:} A) The ability to buy goods and services; B) A financial instrument that represents ownership in a company; C) A system where goods and services are exchanged for commodities; D) A government-issued bond or security; E) A method of payment
\\ \textit{Note:} Credit refers to the ability to access funds or resources to purchase goods and services, which allows consumers or businesses to make transactions without immediate cash payment, thereby facilitating economic activity and growth.
\item \textbf{Answer:} A
\\ \textit{Options:} A) P = ATC \textgreater{} MR = MC; B) P = ATC = MR \textgreater{} MC; C) P \textgreater{} MR = MC = ATC; D) P = MR = MC \textgreater{} ATC; E) P = MR \textgreater{} MC = ATC
\\ \textit{Note:} In the long run, a monopoly firm sets its price (P) above average total cost (ATC) to maximize profits, resulting in price being greater than marginal revenue (MR) and marginal cost (MC), as monopolies face downward-sloping demand curves.
\item \textbf{Answer:} C
\\ \textit{Options:} A) The price of the land is entirely made up of transfer earnings.; B) Economic rent remains constant while transfer earnings fluctuate with market demands.; C) Transfer earnings are zero and the price of the land is entirely made up of economic rent.; D) Both transfer earnings and economic rent increase proportionally.; E) Transfer earnings and economic rent vary inversely; as one increases, the other decreases.
\\ \textit{Note:} When land has only a single use, it means that its value is derived solely from the economic rent it generates, as transfer earnings (the opportunity cost of using the land for its next best alternative) are zero due to the absence of alternative uses.
\item \textbf{Answer:} B
\\ \textit{Options:} A) It is often used to meet emergencies.; B) Its quantity is controlled by the government rather than by market forces.; C) It is backed by a physical commodity like gold or silver.; D) It encourages the accumulation of debt through interest-free loans.; E) It is used for buying military supplies during a war.
\\ \textit{Note:} The danger inherent in the issuance of fiat money lies in the fact that its quantity is subject to government control, which can lead to inflation or hyperinflation if mismanaged, as it is not regulated by the natural supply and demand dynamics of the market.
\end{enumerate}

\section*{True / False}
\begin{enumerate}[label=\arabic*.]
\setcounter{enumi}{5}
\item \textbf{Answer:} False
\\ \textit{Options:} A) True; B) False
\\ \textit{Note:} The statement is false because an increase in demand for exports typically strengthens the dollar as foreign buyers need to purchase more dollars to pay for those exports, while diminishing preferences for foreign goods further support domestic currency value.
\item \textbf{Answer:} False
\\ \textit{Options:} A) True; B) False
\\ \textit{Note:} Capitalism, by its nature, promotes freedom of choice as it encourages competition and innovation, leading to a diverse array of products and services for consumers, rather than restricting options through regulations.
\item \textbf{Answer:} True
\\ \textit{Options:} A) True; B) False
\\ \textit{Note:} A monopsony has significant market power to set wages lower than in a competitive market because it is the sole buyer of labor, allowing it to hire more workers at a reduced price compared to multiple competing employers.
\item \textbf{Answer:} True
\\ \textit{Options:} A) True; B) False
\\ \textit{Note:} In a transition from a monopsony to a perfectly competitive labor market, wages and employment would increase because competition among employers drives wages up and encourages more hiring, as firms are no longer the sole buyers of labor.
\item \textbf{Answer:} False
\\ \textit{Options:} A) True; B) False
\\ \textit{Note:} Economic policy implementation typically involves significant time-lags due to the delays in recognizing economic conditions, formulating appropriate responses, and the time it takes for these policies to take effect and influence the economy.
\end{enumerate}

\section*{Long Form Response}
\begin{enumerate}[label=\arabic*.]
\setcounter{enumi}{10}
\item \textbf{Answer:} The purchase of a Chinese-manufactured entertainment system by an American consumer negatively impacts U.S. national income accounts, specifically through a decline in net exports and GDP. When the American consumer buys the entertainment system, it counts as an import, which increases the value of imports in the national accounts. Since net exports are calculated as the difference between exports and imports, this transaction leads to a decrease in net exports. Consequently, a reduction in net exports directly contributes to a lower Gross Domestic Product (GDP), as GDP accounts for total economic output, including consumption, investment, government spending, and net exports. Thus, both net exports and GDP fall as a result of this purchase.
\\ \textit{Note:} correct because purchasing a Chinese-manufactured entertainment system increases U.S. imports, which reduces net exports and ultimately leads to a decrease in GDP, reflecting the negative impact on U.S. national income accounts.
\item \textbf{Answer:} Monopoly and monopsony, while often confused, represent distinct market structures with significant implications for market dynamics and consumer welfare. A monopoly occurs when a single seller dominates the market, typically leading to higher prices and reduced output, negatively impacting consumer welfare. In contrast, a monopsony arises when a single buyer controls the market, which can suppress prices paid to suppliers, potentially harming producers and limiting the variety of goods available to consumers. Although both structures can lead to inefficiencies and reduced competition, the fundamental difference lies in the direction of market power: monopoly exerts influence over sellers, while monopsony affects buyers. Understanding these differences is crucial for analyzing economic behavior and the overall health of markets in an economy.
\\ \textit{Note:} The answer correctly distinguishes between monopoly and monopsony by highlighting their unique market power dynamics-monopoly leads to higher prices and reduced output harming consumers, while monopsony suppresses supplier prices affecting producers and limiting goods variety-thereby addressing their implications on market dynamics and consumer welfare.
\end{enumerate}

\vfill
\noindent\textit{End of Answer Key}
\end{document}